\section{Provenance, reproduction and data handling}
\label{app:provenance}

\paragraph{Source and build.}
This memo was built from the repository \url{https://github.com/brando90/judges-dont-need-psychometrics-or-corr} on 11-25-2026.
Every experiment result it reads was committed at or before commit \texttt{2bd73d9}, and the memo's own files are in the commit that carries this PDF.
They sit in \texttt{paper\_latex\_and\_notes/ICML\_format\_tech\_memo\_11-25-2026/}.
Running \texttt{python3 -I make\_memo\_assets.py} there regenerates every figure and table of this memo from committed result files and writes \texttt{assets\_log.txt}; \texttt{make} then runs \texttt{pdflatex} and \texttt{bibtex}.
Neither step calls a model.

\paragraph{Where each number comes from.}
We typed each number in the text from a result file or script output committed under \texttt{experiments/}, or from the log of \texttt{make\_memo\_assets.py}, and \texttt{check\_numbers.py} compares them with those sources.
The checker is a screening aid: it accepts a number that appears anywhere in the sources and does not confirm what the number refers to.
\Cref{tab:index} names the folder of each experiment; the folders hold the frozen protocol, the run receipts, request hashes and a reading guide that states which data, metric and comparison each number uses.

\paragraph{Data.}
Public sets (ChaosNLI, DICES, HelpSteer2, KonIQ-10k, SummEval, USR, MovieLens-100K) are used as released.
The retrospective theorem-judge studies use 75 artifacts, 25 tasks and 195 expert ratings, and the homework studies use 160 short answers with two graders; this memo reports only aggregate results from them and reproduces no rater identity and no raw rating.
No new human ratings were collected for the memo.

\paragraph{Prior exposure and the human floor.}
Results on the 75 theorem artifacts are not fresh: experiments 02, 09, 10, 12, 14, 16, 18, 19 and 20 reuse them (Section~\ref{sec:limits}).
The per-item vote counts that a human leave-one-vote-out floor needs for the frozen-judge panels are not saved in the repository, so the memo estimates no such floor; it quotes the reference scores that experiment 29 saved for its test sets as descriptive values.

\paragraph{Cost.}
The scaling-law analysis of experiment 31, the figures and the tables of this memo are free CPU computations on saved results.
The model calls of the underlying experiments (frontier-judge queries through the Stanford AI API Gateway, LoRA training of Qwen3-4B-Instruct-2507 and Llama-3.1-8B-Instruct) are described in the folders of experiments 21, 26, 28 and 29.

\paragraph{Related deposit.}
The identification results of Section~\ref{sec:identification} were deposited as a frozen preprint in the Stanford Digital Repository~\citep{miranda2026correlation}; this memo extends them with proofs, controlled studies, matched-budget comparisons and scaling fits.
